\documentclass[11pt,a4paper]{article}

% Same restrained layout as the companion preprint; portable under tectonic.
\usepackage[top=2.65cm,bottom=2.7cm,left=3.0cm,right=3.0cm]{geometry}
\usepackage[T1]{fontenc}
\usepackage{lmodern}
\usepackage{amsmath,amssymb,amsthm,mathtools}
\usepackage{booktabs}
\usepackage{enumitem}
\usepackage{fancyhdr}
\usepackage{microtype}
\usepackage{xcolor}
\usepackage[colorlinks=true,linkcolor=blue!55!black,citecolor=blue!55!black,urlcolor=blue!55!black]{hyperref}
\usepackage{bookmark}

\setlength{\parindent}{1.25em}
\setlength{\parskip}{0.25em}
\setlength{\emergencystretch}{2em}
\setlist[itemize]{leftmargin=2em,itemsep=0.15em,topsep=0.25em}
\setlist[enumerate]{leftmargin=2em,itemsep=0.15em,topsep=0.25em}
\setcounter{tocdepth}{2}

\pagestyle{fancy}
\fancyhf{}
\fancyhead[L]{\small\itshape Infinite Fourier tails of the semilocal Weil form}
\fancyhead[R]{\small\thepage}
\renewcommand{\headrulewidth}{0.35pt}
\fancypagestyle{plain}{%
  \fancyhf{}%
  \fancyfoot[C]{\small\thepage}%
  \renewcommand{\headrulewidth}{0pt}%
}

\newtheorem{theorem}{Theorem}[section]
\newtheorem{lemma}[theorem]{Lemma}
\newtheorem{corollary}[theorem]{Corollary}
\newtheorem{proposition}[theorem]{Proposition}
\theoremstyle{definition}
\newtheorem{definition}[theorem]{Definition}
\newtheorem{remark}[theorem]{Remark}
\newtheorem{problem}[theorem]{Open Problem}

\newtheorem*{mainthm}{Main Theorem}

\newcommand{\RH}{\mathrm{RH}}
\newcommand{\Z}{\mathbb{Z}}
\newcommand{\R}{\mathbb{R}}
\newcommand{\C}{\mathbb{C}}
\newcommand{\N}{\mathbb{N}}
\newcommand{\eps}{\varepsilon}
\newcommand{\ip}[2]{\langle #1,#2\rangle}
\newcommand{\norm}[1]{\lVert #1\rVert}
\newcommand{\abs}[1]{\lvert #1\rvert}
\newcommand{\Lmin}{L_{-}}
\newcommand{\Lmax}{L_{+}}
\newcommand{\Arch}{\mathfrak{A}}
\newcommand{\Dist}{\mathfrak{D}}
\newcommand{\sinw}{\mathrm{s}}
\DeclareMathOperator{\Ima}{Im}
\DeclareMathOperator{\Rea}{Re}
\DeclareMathOperator{\diag}{diag}

\title{\bfseries\Large Infinite Fourier Tails and High-Mode Coercivity\\ for the Semilocal Weil Quadratic Form\\[0.35em]
\large A complete proof of a coupled tail bound for the rank-eight ground state\\ on the support cell $17\le C\le 19$}

\author{Biswajit Mondal}
\date{\small October 6, 2026 \quad doi:10.5281/zenodo.23187645\\[0.25em]
\textit{Research preprint. The Riemann Hypothesis is not claimed.}}

\begin{document}
\maketitle

\begin{abstract}
Connes, Consani and Moscovici encode the semilocal Weil quadratic form on
$[C^{-1/2},C^{1/2}]$ by an infinite real symmetric matrix $A(C)$ in a
Fourier basis, with off-diagonal entries of displacement form
$A_{ij}=(b_i-b_j)/(i-j)$. Their spectral programme compares the ground
states of finite truncations of $A(C)$ with the zeros of the Riemann zeta
function. Any such comparison has to control how the infinite matrix acts on
a truncated ground state. We derive the matrix from the explicit-formula
distributions and prove three rank-independent estimates. First, an exact
identity for the omitted components of $A\,c$, for any finitely supported
even vector $c$, expressed through signed moments of $c$ plus an explicit
remainder that is uniform over infinitely many indices. Second, the operator
bound $\norm{A_{\mathrm{off}}}\le 2\pi\sup_i\abs{b_i}$, obtained from the
discrete Hilbert transform. Third, explicit logarithmic coercivity of
high-mode compressions. We combine these with validated interval
enclosures of the rank-eight even ground vector $u(C)$. The result is the
following uniform statement over the whole cell $17\le C\le 19$: the
even components of $A(C)\,u(C)$ at every Fourier index $j>64$ have joint
$\ell^2$ norm below $9.491\times 10^{-12}$. Ignoring the signs gives a
bound of about $2.906$. Moreover, the even compression of $A(C)$ to
indices $j\ge 18{,}612{,}929$, and the full compression to
$\abs{j}\ge 60{,}879{,}750$, both exceed the identity. Every
infinite-index statement is proved analytically; finite-dimensional
quantities are certified by outward-rounded ball arithmetic and replayed
by an independent implementation. These results do not bound the response
of the intermediate modes $9\le j<18{,}612{,}929$. They give no uniform
profile budget over all ranks and supports, and they do not identify any
limit with Riemann's $\Xi$ function. We do not prove the Riemann Hypothesis,
and we explain why the present estimates cannot by themselves imply it.
\end{abstract}

\tableofcontents

%======================================================================
\section{Introduction}\label{sec:intro}
%======================================================================

\subsection{Context}
Weil's explicit formula expresses a sum over the nontrivial zeros of
$\zeta$ as a distribution $\Psi$ built from the pole of $\zeta$, the
archimedean place and the primes. Weil's criterion~\cite{Weil1952} states
that the Riemann Hypothesis holds if and only if the quadratic form
$f\mapsto \Psi(f^{*}*f)$ is nonnegative on a suitable class of test functions;
see also Bombieri~\cite{Bombieri2000}. Restricting the test functions to
$[\lambda^{-1},\lambda]\subset\R_+^{*}$ produces the \emph{semilocal} form
$QW_\lambda$. Only the finitely many prime powers $p^a\le\lambda^2$ enter it.

Connes, Consani and Moscovici~\cite{CCM2025} computed the matrix of
$QW_\lambda$ in the Fourier basis of $L^2([\lambda^{-1},\lambda],d^{*}u)$.
They observed that the off-diagonal part has displacement form, and
constructed from the even ground state of each finite truncation a
self-adjoint operator. The spectrum of that operator consists of the real
zeros of an entire function. Numerically these zeros approximate the
lowest zeros of $\zeta(\tfrac12+is)$ with striking accuracy. Proving
convergence would prove RH~\cite[\S7--8]{CCM2025}.

A necessary ingredient of any convergence argument is control of the
\emph{infinite} matrix acting on a \emph{truncated} ground state. If
$c$ is the ground vector at retained Fourier rank $N$, then $A c$ has
infinitely many nonzero components outside the retained block. The
retained eigenproblem cannot see these components. A naive row bound on
them fails: individual off-diagonal rows decay only like $1/\abs{i-j}$ and
are not absolutely summable. This paper proves an estimate of this kind
that is uniform over all omitted indices and over a whole support interval.

\subsection{Main results}
Throughout, $C>1$ is the support parameter ($C=\lambda^2$ in the notation of
\cite{CCM2025}), $L=\log C$, and $A=A(C)$ is the infinite Weil matrix of
Definition~\ref{def:matrix}. The even orthonormal coordinates of a vector are
$u=(u_0,\dots,u_N)$, with full coefficients $c_0=u_0$ and
$c_{\pm i}=u_i/\sqrt2$. For $n>N$ the even component of $Ac$ at index $n$ is
$R_n=\sqrt2\,(Ac)_n$.

\begin{mainthm}
Let $17\le C\le 19$, and let $u(C)$ be the unit ground eigenvector, oriented
by $u_0>0$, of the even block of the $17\times17$ truncation of $A(C)$
(retained rank $N=8$). Then:
\begin{enumerate}[label=\textup{(\roman*)}]
\item \textbf{Infinite forcing tail.}\quad
$\displaystyle\Bigl(\sum_{n>64}R_n(C)^2\Bigr)^{1/2}<9.491\times10^{-12}$
for every $C\in[17,19]$. (Theorem~\ref{thm:main-tail}.)
\item \textbf{Rank-independent off-diagonal bound.}\quad The off-diagonal
part of $A(C)$ defines a bounded operator on $\ell^2(\Z)$ with norm below
$19.836078$. Every finite compression obeys the same bound.
(Theorem~\ref{thm:hilbert}, Proposition~\ref{prop:b-bounds}.)
\item \textbf{High-mode coercivity.}\quad For every nonzero finitely supported
$x$ with $x_j=0$ for $\abs{j}<n_0$:
$\ip{x}{A(C)x}>\norm{x}^2$ if $x$ is even and $n_0\ge18{,}612{,}929$; and
$\ip{x}{A(C)x}>\norm{x}^2$ for arbitrary $x$ if $n_0\ge60{,}879{,}750$.
(Theorem~\ref{thm:coercive}.)
\end{enumerate}
\end{mainthm}

The infinite range of indices in (i) and the rank independence in (ii)
and (iii) are proved analytically (Sections~\ref{sec:tail}--\ref{sec:coercive}).
The only computer-assisted inputs are enclosures of finitely many real
numbers: the boundary sum, sixteen signed moments and two absolute moments
of $u(C)$, uniformly over the cell. Section~\ref{sec:certificate} gives
their validation theorems in full.

The decisive point is \emph{order of operations}. The leading coefficient
of $R_n$ is $\sqrt2\,b_n S/n$, where $S=\sum_j c_j$ is the boundary value of
the ground state. For the ground state $\abs{S}<6.6\times10^{-12}$
(Table~\ref{tab:worst-cell}). If $S$ is formed before absolute values are
taken, this cancellation survives. Expanding the remaining terms in signed
moments preserves further cancellation. If instead each entry is bounded in
absolute value and summed, the bound is $2.906$, about $3\times10^{11}$
times larger (Remark~\ref{rem:uncoupled}).

\subsection{What is not proved}\label{subsec:not-proved}
We state the limits of these results plainly, because the surrounding
programme aims at RH.
\begin{itemize}
\item The theorems control the forcing at indices $j>64$ and the far block
$j\ge 18{,}612{,}929$. They do \emph{not} control the inverse response of
the intermediate modes $9\le j<18{,}612{,}929$. Eliminating those modes can
amplify the forcing through a nearly singular constrained inverse:
the certified ground eigenvalue lies in
$[2.26\times10^{-26},1.38\times10^{-24}]$, and the spectral gap is about
$4\times10^{-23}$ (Section~\ref{sec:limits}).
\item No estimate here is uniform in growing support $C\to\infty$, and none
provides a summable support/rank path.
\item No limit is identified with Riemann's $\Xi$ function, and
the Riemann Hypothesis is not proved.
\end{itemize}

\subsection{Organization and conventions}
Section~\ref{sec:matrix} derives the matrix from the explicit-formula
distributions. Section~\ref{sec:closed} evaluates its entries in closed form.
Section~\ref{sec:operator} proves the uniform displacement bounds and the
Hilbert-commutator estimate. Section~\ref{sec:tail} proves the omitted-mode
identity and its infinite-index moment bound. Section~\ref{sec:coercive}
proves high-mode coercivity. Section~\ref{sec:certificate} proves the
validation lemmas and states the computer-assisted theorem.
Section~\ref{sec:limits} records a consequence and the remaining
obstructions. Appendix~\ref{app:crosscheck} reports an independent
cross-check, and Appendix~\ref{app:repro} gives replay instructions and
hashes.

We use the following classical facts without proof: the Weierstrass
product of $\Gamma$ in the form~\eqref{eq:psi-series}, Gauss's integral for
$\psi$, the expansion $E_1(x)=-\gamma-\log x+O(x)$ of the exponential
integral~\cite[5.7.6, 5.9.16, 6.6.2]{DLMF}, Parseval's identity on the
circle, and the inclusion property of interval arithmetic~\cite{Johansson2017,Rump2010}.
Everything else is proved here.

%======================================================================
\section{The semilocal Weil form and its matrix}\label{sec:matrix}
%======================================================================

\subsection{The explicit-formula distributions}
For a function $F$ on $\R_+^{*}$ we write $F^{\vee}(x)=F(x^{-1})$ and
$d^{*}x=dx/x$. Following~\cite{Weil1952,CCM2025}, the distributions of the
explicit formula in the normalization $F(x)=x^{1/2}f(x)$ are
\begin{align}
W_{0,2}(F)&=\int_0^\infty F(x)\,(x^{1/2}+x^{-1/2})\,d^{*}x,\label{eq:W02}\\
W_p(F)&=(\log p)\sum_{m\ge1}p^{-m/2}\bigl(F(p^m)+F(p^{-m})\bigr),\label{eq:Wp}\\
W_\R(F)&=(\log4\pi+\gamma)F(1)+\int_1^\infty\Bigl(F(x)+F(x^{-1})-2x^{-1/2}F(1)\Bigr)\frac{x^{1/2}}{x-x^{-1}}\,d^{*}x.\label{eq:WR}
\end{align}
Put $\Psi=W_{0,2}-W_\R-\sum_pW_p$. With the involutive convolution algebra
$(f^{*}*g)(y)=\int\overline{f(x/y)}\,g(x)\,d^{*}x$ on $\R_+^{*}$, the Weil
sesquilinear form is $QW(f,g)=\Psi(f^{*}*g)$.

\begin{lemma}[Half-line reduction]\label{lem:halfline}
Define on functions on $[1,\infty)$
\[
W_{0,2}^{\sharp}(F)=\int_1^\infty F(x)(x^{1/2}+x^{-1/2})\,d^{*}x,\qquad
W_p^{\sharp}(F)=(\log p)\sum_{m\ge1}p^{-m/2}F(p^m),
\]
\[
W_\R^{\sharp}(F)=\tfrac12(\log4\pi+\gamma)F(1)+\int_1^\infty\frac{x^{1/2}F(x)-F(1)}{x-x^{-1}}\,d^{*}x,
\]
and $\Psi^{\sharp}=W_{0,2}^{\sharp}-W_\R^{\sharp}-\sum_pW_p^{\sharp}$. Then
$\Psi(h)=\Psi^{\sharp}(h+h^{\vee})$ for every $h$ for which both sides are
defined.
\end{lemma}

\begin{proof}
Each identity follows by splitting $\int_0^\infty$ at $x=1$ and substituting
$x\mapsto x^{-1}$ on $(0,1)$. The measure $d^{*}x$ is invariant under this
substitution, and $x^{1/2}+x^{-1/2}$ is symmetric. For $W_\R$ note that
$(h+h^{\vee})(1)=2h(1)$, which accounts for the factor $\tfrac12$ in front of
the constant and for the term $2x^{-1/2}F(1)$ in~\eqref{eq:WR}.
\end{proof}

\subsection{Additive coordinates and the Fourier basis}
Fix $C>1$ and $L=\log C$. The map $u\mapsto \log(C^{1/2}u)$ carries
$[C^{-1/2},C^{1/2}]$ with $d^{*}u$ isometrically onto $[0,L]$ with $dx$. It
carries the multiplicative convolution to the additive one,
$(f^{*}*g)(y)=\int\overline{f(x-y)}\,g(x)\,dx$, and inversion to $y\mapsto -y$.
We work on $\mathcal H_L=L^2([0,L],dx)$ with the orthonormal basis
\begin{equation}\label{eq:basis}
U_n(x)=L^{-1/2}e^{iw_nx},\qquad w_n=\frac{2\pi n}{L},\qquad n\in\Z.
\end{equation}
Set $q(f,g)(y)=(f^{*}*g)(y)+(f^{*}*g)(-y)$, an even function supported in
$[-L,L]$. Pulled back to $[0,L]$ (where $y=\log x$, $x\ge1$), the
half-line distributions become
\begin{equation}\label{eq:Dist}
\Dist(F)=\int_0^LF(y)\,2\cosh\tfrac y2\,dy-\Arch(F)-\sum_{1<p^a<C}\frac{\log p}{p^{a/2}}F(a\log p),
\end{equation}
\begin{equation}\label{eq:Arch}
\Arch(F)=\tfrac12(\log4\pi+\gamma)F(0)+\int_0^\infty\frac{e^{y/2}F(y)-F(0)}{2\sinh y}\,dy,
\end{equation}
for $F$ continuous on $[0,\infty)$, piecewise $C^1$, and vanishing on
$[L,\infty)$. A prime power with $p^a=C$ would contribute $F(L)=0$, so
the sum may equally run over $p^a\le C$. In particular, on
$17\le C\le 19$ the active set is $\{p^a\le17\}$ throughout; at $C=19$ the
term $p^a=19$ vanishes identically.

\begin{lemma}[Correlations]\label{lem:corr}
For $y\in[0,L]$,
\[
q(U_m,U_n)(y)=\begin{cases}\dfrac{\sin(w_my)-\sin(w_ny)}{\pi(n-m)},&m\ne n,\\[2mm]
2\bigl(1-\tfrac yL\bigr)\cos(w_ny),&m=n.\end{cases}
\]
\end{lemma}

\begin{proof}
For $0\le y\le L$, $(U_m^{*}*U_n)(y)=L^{-1}\int_y^Le^{iw_m(y-x)}e^{iw_nx}\,dx$.
If $m\ne n$ this equals
$\bigl(e^{iw_my}-e^{iw_ny}\bigr)/\bigl(2\pi i(n-m)\bigr)$, because
$e^{i(w_n-w_m)L}=1$. If $m=n$ it equals $(1-y/L)e^{iw_ny}$. Since
$(f^{*}*g)(-y)=\overline{(g^{*}*f)(y)}$ and the expression for $m\ne n$ is
symmetric in $(m,n)$, the function $q$ is twice the real part. This gives
the stated formulas.
\end{proof}

\begin{definition}[Weil matrix]\label{def:matrix}
The \emph{Weil matrix} at support $C$ is the infinite real matrix
$A_{mn}=QW(U_m,U_n)=\Dist\bigl(q(U_m,U_n)\bigr)$, for $m,n\in\Z$. Write
$A=P-H-S$ for the pole, archimedean and prime contributions, i.e.\ the three
terms of~\eqref{eq:Dist}. For a finitely supported $c\in\C^{\Z}$ with
$f=\sum_nc_nU_n$, one has $QW(f,f)=\ip{c}{Ac}$.
\end{definition}

\begin{proposition}[Displacement structure]\label{prop:displacement}
Let $\sinw_n(y)=\sin(w_ny)\mathbf 1_{[0,L]}(y)$ and $b_n=-\pi^{-1}\Dist(\sinw_n)$.
Then $b_{-n}=-b_n$, $b_0=0$, $A_{-m,-n}=A_{mn}=A_{nm}$, and
\begin{equation}\label{eq:displacement}
A_{mn}=\frac{b_m-b_n}{m-n}\quad(m\ne n),\qquad b_n=n\,A_{n0}.
\end{equation}
The same holds separately for each of the three terms of~\eqref{eq:Dist},
so that $b=b^P+b^\R+b^\Pi$ with each part odd in $n$.
\end{proposition}

\begin{proof}
$\Dist$ is linear, and $\sinw_n(0)=0$, so~\eqref{eq:Arch} needs no
regularization on $\sinw_n$. By Lemma~\ref{lem:corr},
$A_{mn}=\Dist(\sinw_m-\sinw_n)/(\pi(n-m))=(b_m-b_n)/(m-n)$. Oddness of
$n\mapsto\sinw_n$ gives $b_{-n}=-b_n$ and $b_0=0$. Taking $n=0$ gives
$b_m=mA_{m0}$. The diagonal formula of Lemma~\ref{lem:corr} is even in $n$,
and~\eqref{eq:displacement} is invariant under $(m,n)\mapsto(-m,-n)$. All
three distributions are real on real functions.
\end{proof}

\begin{remark}
Equivalently, with $\mathcal J=\diag(n)$ and $e=(1)_{n}$, one has the commutator
identity $[A,\mathcal J]=e\,b^{T}-b\,e^{T}$ on finitely supported vectors. This is the
structure used in~\cite[Lemma~5.2]{CCM2025}.
\end{remark}

Let $\gamma$ be the reflection $U_n\mapsto U_{-n}$. Under the centring
$x\mapsto x-L/2$ it is the reflection $x\mapsto-x$, which corresponds to
$u\mapsto u^{-1}$ on $[C^{-1/2},C^{1/2}]$. Since $A$ commutes with $\gamma$,
$A$ preserves the even vectors $c_{-n}=c_n$. The even orthonormal basis is
$e_0=U_0$ and $e_j=(U_j+U_{-j})/\sqrt2$ for $j\ge1$.

%======================================================================
\section{Closed formulas for the entries}\label{sec:closed}
%======================================================================

\subsection{Digamma facts}
Let $\psi=\Gamma'/\Gamma$. We use
\begin{equation}\label{eq:psi-series}
\psi(z)=-\gamma+\sum_{k\ge0}\Bigl(\frac1{k+1}-\frac1{k+z}\Bigr),\qquad
\psi'(z)=\sum_{k\ge0}\frac1{(k+z)^2},
\end{equation}
\begin{equation}\label{eq:gauss}
\psi(z)=\int_0^\infty\Bigl(\frac{e^{-t}}t-\frac{e^{-zt}}{1-e^{-t}}\Bigr)dt,\qquad
\psi'(z)=\int_0^\infty\frac{t\,e^{-zt}}{1-e^{-t}}\,dt\qquad(\Rea z>0).
\end{equation}
The second identity in~\eqref{eq:gauss} follows from~\eqref{eq:psi-series}
by writing $(k+z)^{-2}=\int_0^\infty te^{-(k+z)t}dt$ and summing the
geometric series, which Tonelli's theorem permits. From the first identity
in~\eqref{eq:gauss} and $E_1(x)=\int_x^\infty e^{-t}t^{-1}dt=-\gamma-\log x+O(x)$,
\begin{equation}\label{eq:reg}
\int_\eps^\infty\frac{e^{-zt}}{1-e^{-t}}\,dt=-\gamma-\log\eps-\psi(z)+o(1)\qquad(\eps\downarrow0).
\end{equation}
Let $a_k=2k+\tfrac12$ and $\rho(y)=e^{y/2}/(2\sinh y)$. Then
\begin{equation}\label{eq:rho}
\rho(y)=\sum_{k\ge0}e^{-a_ky}\quad(y>0),\qquad
\int_\eps^\infty\rho(y)e^{iwy}dy=\tfrac12\int_{2\eps}^\infty\frac{e^{-zt}}{1-e^{-t}}dt,\quad z=\tfrac14-\tfrac{iw}2,
\end{equation}
by the substitution $t=2y$.

\subsection{The three contributions}
Write $D=\cosh(L/2)-1$ and $z_n=\tfrac14+i\pi n/L=\tfrac14+iw_n/2$.

\begin{proposition}[Pole]\label{prop:pole}
For $m\ne n$,
$P_{mn}=\dfrac{2D}{\pi(m-n)}\Bigl[\dfrac{w_m}{\frac14+w_m^2}-\dfrac{w_n}{\frac14+w_n^2}\Bigr]$,
and $P_{nn}=\dfrac{4D}{L}\cdot\dfrac{\frac14-w_n^2}{(\frac14+w_n^2)^2}$.
\end{proposition}

\begin{proof}
For $s=a+iw_n$ with $a=\pm\frac12$ we have $e^{sL}=e^{aL}$, hence
$\int_0^L e^{sy}\sin(w_n y)\,dy$ and the corresponding integral with
$(1-y/L)$ are elementary:
$\int_0^L e^{ay}\sin(w_ny)\,dy=w_n(1-e^{aL})/(a^2+w_n^2)$ and
$\int_0^L(1-\tfrac yL)e^{sy}dy=-\tfrac1s+\tfrac{e^{aL}-1}{Ls^2}$. Summing
over $a=\pm\frac12$ with $2\cosh\frac y2=e^{y/2}+e^{-y/2}$ and taking real
parts (the terms $-1/s$ cancel in pairs) gives
$\int_0^L\sinw_n\,2\cosh\frac y2\,dy=-2Dw_n/(\frac14+w_n^2)$ and the stated
diagonal. The off-diagonal formula follows from
Proposition~\ref{prop:displacement}.
\end{proof}

\begin{proposition}[Archimedean place]\label{prop:arch}
For every $n\in\Z$,
\[
H_{nn}=\log\pi-\Rea\psi(z_n)-\frac{\Rea\psi'(z_n)}{2L}+\frac2L\sum_{k\ge0}e^{-a_kL}\frac{a_k^2-w_n^2}{(a_k^2+w_n^2)^2},
\]
and $\Arch(\sinw_n)=\tfrac12\Ima\psi(z_n)-\sum_{k\ge0}e^{-a_kL}\dfrac{w_n}{a_k^2+w_n^2}$.
\end{proposition}

\begin{proof}
Let $F=2(1-y/L)\cos(w_ny)\mathbf 1_{[0,L]}$, so $F(0)=2$. The integrand of
\eqref{eq:Arch} is bounded near $0$ and integrable at $\infty$, so
\[
\Arch(F)=\log4\pi+\gamma+\lim_{\eps\downarrow0}\Bigl[\int_\eps^L\rho F\,dy-\int_\eps^\infty\frac{dy}{\sinh y}\Bigr].
\]
First, $\int_\eps^\infty dy/\sinh y=\log\coth(\eps/2)=-\log(\eps/2)+O(\eps^2)$.
Second, split $\int_\eps^L\rho F$ as
$2\Rea\int_\eps^\infty\rho e^{iw_ny}-\tfrac2L\Rea\int_\eps^\infty y\rho e^{iw_ny}+2\int_L^\infty\rho(\tfrac yL-1)\cos(w_ny)$.
By~\eqref{eq:rho} and~\eqref{eq:reg} the first piece is
$-\gamma-\log(2\eps)-\Rea\psi(\bar z_n)+o(1)$. By~\eqref{eq:gauss} with $t=2y$
the second piece tends to $-\Rea\psi'(\bar z_n)/(2L)$. In the third piece put
$y=L+s$ and use $\cos(w_n(L+s))=\cos(w_ns)$ together with~\eqref{eq:rho}:
it equals
$\tfrac2L\sum_ke^{-a_kL}\Rea(a_k-iw_n)^{-2}=\tfrac2L\sum_ke^{-a_kL}(a_k^2-w_n^2)/(a_k^2+w_n^2)^2$.
The constants combine to
$\log4\pi+\gamma+\log(\eps/2)-\gamma-\log(2\eps)=\log\pi$. Finally
$\Rea\psi(\bar z)=\Rea\psi(z)$ and $\Rea\psi'(\bar z)=\Rea\psi'(z)$, and
$H_{nn}=\Arch(F)$.

For $\sinw_n$ we have $F(0)=0$, so $\Arch(\sinw_n)=\int_0^L\rho\,\sinw_n$.
The imaginary part of~\eqref{eq:reg} has a limit, and it gives
$\Ima\int_0^\infty\rho e^{iw_ny}dy=-\tfrac12\Ima\psi(\bar z_n)=\tfrac12\Ima\psi(z_n)$.
The tail $\int_L^\infty\rho\sin(w_ny)\,dy$ equals
$\sum_ke^{-a_kL}w_n/(a_k^2+w_n^2)$, by the same substitution.
\end{proof}

\begin{proposition}[Primes]\label{prop:primes}
$S_{nn}=\sum_{p^a<C}\frac{\log p}{p^{a/2}}\,2\bigl(1-\frac{a\log p}{L}\bigr)\cos(w_na\log p)$,
and the prime contribution to $b_n$ is
$b^\Pi_n=\pi^{-1}\sum_{p^a<C}\frac{\log p}{p^{a/2}}\sin(w_na\log p)$.
\end{proposition}

\begin{proof}
Evaluate the finite sum in~\eqref{eq:Dist} on Lemma~\ref{lem:corr}.
\end{proof}

\begin{corollary}[Closed displacement sequence]\label{cor:b}
For $n\ge1$,
\begin{equation}\label{eq:b-closed}
b_n=\frac{2D}{\pi}\frac{w_n}{\frac14+w_n^2}
+\underbrace{\frac{\Ima\psi(z_n)}{2\pi}-\frac1\pi\sum_{k\ge0}e^{-a_kL}\frac{w_n}{a_k^2+w_n^2}}_{=:\,b_n^\R}
+\frac1\pi\sum_{p^a<C}\frac{\log p}{p^{a/2}}\sin(w_na\log p).
\end{equation}
\end{corollary}

\begin{proof}
Use $b_n=-\pi^{-1}\Dist(\sinw_n)$ together with
Propositions~\ref{prop:pole}--\ref{prop:primes}. The archimedean
distribution enters $\Dist$ with a minus sign, so it contributes
$+\pi^{-1}\Arch(\sinw_n)$.
\end{proof}

\begin{remark}
These closed forms are equivalent to the hypergeometric expressions
of~\cite[Prop.~4.2--4.3]{CCM2025}. The constant $\log\pi$ is obtained from
the regularization, not fitted. Appendix~\ref{app:crosscheck} compares the
closed entries with direct numerical quadrature of~\eqref{eq:Dist}--\eqref{eq:Arch},
using code that is independent of the certificate generator. They agree
to better than $2\times10^{-30}$.
\end{remark}

%======================================================================
\section{Rank-independent operator bounds}\label{sec:operator}
%======================================================================

From now on $17\le C\le19$. Put $\Lmin=\log17$, $\Lmax=\log19$, and
\begin{equation}\label{eq:constants}
E=\frac{e^{-\Lmin/2}}{1-e^{-2\Lmin}},\qquad
\mathcal P=\sum_{p^a\le17}\frac{\log p}{\sqrt{p^a}},\qquad
D_{+}=\cosh\tfrac{\Lmax}2-1.
\end{equation}
Since $e^{\Lmin/2}=\sqrt{17}$ and $e^{\Lmax/2}=\sqrt{19}$, these are explicit
algebraic and logarithmic numbers. Table~\ref{tab:constants} gives outward
enclosures.

\begin{lemma}[Digamma inequalities]\label{lem:digamma}
For $y>0$:
\begin{enumerate}[label=\textup{(\alph*)}]
\item $0<\Ima\psi(\tfrac14+iy)\le\tfrac\pi2+\min(2,1/y)$;
\item $\Rea\psi(\tfrac14+iy)\ge\log y-1/y$;
\item $\abs{\psi'(\tfrac14+iy)}\le 1/y^2+\pi/(2y)$.
\end{enumerate}
\end{lemma}

\begin{proof}
(a) By~\eqref{eq:psi-series},
$\Ima\psi(\frac14+iy)=\sum_{k\ge0}g(k)$ with $g(t)=y/((t+\frac14)^2+y^2)>0$
decreasing on $[0,\infty)$. Hence the sum is at most
$g(0)+\int_0^\infty g=g(0)+\frac\pi2-\arctan\frac1{4y}$. Finally
$g(0)=y/(\frac1{16}+y^2)\le\min(2,1/y)$, since the maximum $2$ is attained
at $y=\frac14$.

(b) By~\eqref{eq:psi-series} and $\gamma=\lim_M(\sum_{k<M}\frac1{k+1}-\log M)$,
we have $\Rea\psi(\frac14+iy)=\lim_M[\log M-\sum_{k<M}f(k)]$ with
$f(t)=(t+\frac14)/((t+\frac14)^2+y^2)$. For each $k$,
$\abs{f(k)-\int_k^{k+1}f}\le\int_k^{k+1}\abs{f'}$. The function $f$ increases
up to $t+\frac14=y$ and then decreases, with maximum $1/(2y)$. Hence
$\int_0^\infty\abs{f'}\le 2\cdot\frac1{2y}=\frac1y$, and so
$\sum_{k<M}f(k)\le\int_0^Mf+\frac1y$. Since
$\int_0^Mf=\frac12\log\frac{(M+1/4)^2+y^2}{1/16+y^2}$, letting $M\to\infty$
gives $\Rea\psi\ge\frac12\log(\frac1{16}+y^2)-\frac1y\ge\log y-\frac1y$.

(c) By~\eqref{eq:psi-series},
$\abs{\psi'}\le\sum_k((k+\frac14)^2+y^2)^{-1}$. The $k=0$ term is at most
$y^{-2}$. The sum of the rest is at most
$\int_0^\infty dt/((t+\frac14)^2+y^2)\le\pi/(2y)$.
\end{proof}

\begin{proposition}[Uniform displacement bounds]\label{prop:b-bounds}
For all $C\in[17,19]$ and all $n\in\Z$,
\begin{align}
\abs{b_n}&\le B:=\frac{2D_{+}}\pi+\frac{\frac\pi2+\min(2,\Lmax/\pi)}{2\pi}+\frac E\pi+\frac{\mathcal P}\pi,\label{eq:B}\\
\abs{b^\R_n}&\le B_\R:=\frac{\frac\pi2+\min(2,\Lmax/\pi)}{2\pi}+\frac E\pi,\label{eq:BR}\\
\abs{b_n}&\le B_0+\frac{B_1}{\abs n}\ (n\ne0),\qquad B_0=\frac14+\frac{\mathcal P}\pi,\quad B_1=\frac{\Lmax}{\pi^2}\Bigl(D_{+}+\frac{1+E}2\Bigr).\label{eq:B01}
\end{align}
\end{proposition}

\begin{proof}
By oddness it suffices to take $n\ge1$; let $y=\pi n/L\ge\pi/\Lmax$. In
\eqref{eq:b-closed}: the pole term is at most $2D_{+}/\pi$, because
$w/(\frac14+w^2)\le1$, and also at most $2D_{+}/(\pi w_n)=D_{+}L/(\pi^2n)$.
The digamma term is controlled by Lemma~\ref{lem:digamma}(a), using
$1/y\le\Lmax/\pi$ for~\eqref{eq:B} and $1/y=L/(\pi n)$ for~\eqref{eq:B01}.
For the exponential series, $w/(a^2+w^2)\le1/(2a)\le1$ and also
$\le1/w$; moreover $\sum_ke^{-a_kL}=e^{-L/2}/(1-e^{-2L})\le E$, because this
function decreases in $L$. The prime sum has absolute value at most
$\mathcal P/\pi$, because the active set is $\{p^a\le17\}$ on the whole
cell. Adding the pieces gives the three bounds.
\end{proof}

\begin{theorem}[Hilbert commutator bound]\label{thm:hilbert}
Let $\beta\in\ell^\infty(\Z)$ and let $T_{ij}=(\beta_i-\beta_j)/(i-j)$ for
$i\ne j$, with $T_{ii}=0$. Then $T$ extends from finitely supported vectors
to a bounded operator on $\ell^2(\Z)$ with
$\norm T\le2\pi\norm\beta_\infty$. For every set $I\subset\Z$, the
compression $\Pi_IT\Pi_I$ obeys the same bound.
\end{theorem}

\begin{proof}
Let $\mathsf H_{ij}=1/(i-j)$ for $i\ne j$ and $\mathsf H_{ii}=0$. On
$(0,2\pi)$ the function $\pi-\theta$ has Fourier coefficients
$\frac1{2\pi}\int_0^{2\pi}(\pi-\theta)e^{-ik\theta}d\theta=-i/k$ for $k\ne0$,
and $0$ for $k=0$, by one integration by parts. Hence
$h(\theta):=\sum_{k\ne0}e^{ik\theta}/k=i(\pi-\theta)$ in $L^2(0,2\pi)$. For
finitely supported $x,y$ put $X(\theta)=\sum x_je^{ij\theta}$ and
$Y(\theta)=\sum y_je^{ij\theta}$. Since $\overline XY$ is a trigonometric
polynomial,
\[
\ip x{\mathsf Hy}=\sum_{i\ne j}\frac{\overline{x_i}y_j}{i-j}=\frac1{2\pi}\int_0^{2\pi}h\,\overline XY\,d\theta,
\]
so $\abs{\ip x{\mathsf Hy}}\le\pi\cdot\frac1{2\pi}\int\abs X\abs Y\le\pi\norm x\norm y$
by Cauchy--Schwarz and Parseval. Thus $\norm{\mathsf H}\le\pi$; this is Hilbert's inequality~\cite[Ch.~IX]{HLP}. Now
$T=D_\beta\mathsf H-\mathsf HD_\beta$ with $D_\beta=\diag(\beta_i)$, so
$\norm T\le2\norm\beta_\infty\norm{\mathsf H}$. The compression is
$D_{\Pi\beta}(\Pi\mathsf H\Pi)-(\Pi\mathsf H\Pi)D_{\Pi\beta}$, and
$\norm{\Pi\mathsf H\Pi}\le\pi$.
\end{proof}

\begin{corollary}\label{cor:offdiag}
On the whole cell, the off-diagonal part of $A(C)$ has norm at most
$2\pi B<19.836078$. The off-diagonal part of the archimedean contribution
has norm at most $2\pi B_\R<2.994796$.
\end{corollary}

Individual rows of $A_{\mathrm{off}}$ are not absolutely summable, because
$\sum_j\abs{b_i-b_j}/\abs{i-j}$ diverges logarithmically in general.
Theorem~\ref{thm:hilbert} retains the signed structure that an
absolute-row (Schur) bound discards.

\begin{table}[ht]
\centering
\caption{Constants for the cell $17\le C\le19$. Each enclosure is outward;
the certificate (1024-bit Arb) and the independent check (mpmath interval
arithmetic, 40 digits) agree.}\label{tab:constants}
\begin{tabular}{@{}lll@{}}
\toprule
Quantity & Definition & Enclosure\\
\midrule
$\mathcal P$ & \eqref{eq:constants} & $5.8323263624721337\ldots$\\
$E$ & \eqref{eq:constants} & $0.2433777626232646\ldots$\\
$B$ & \eqref{eq:B} & $<3.157010$\\
$B_{64}=B_0+B_1/64$ & \eqref{eq:B01} & $<2.115418$\\
$2\pi B$ & Cor.~\ref{cor:offdiag} & $<19.836078$\\
$2\pi B_\R$ & Cor.~\ref{cor:offdiag} & $<2.994796$\\
$2\sinh(\Lmax/2)-\Lmax$ & Lemma~\ref{lem:pole-form} & $<1.185045$\\
\bottomrule
\end{tabular}
\end{table}

%======================================================================
\section{The omitted-mode identity and its infinite moment bound}\label{sec:tail}
%======================================================================

Fix a retained rank $N\ge1$ and real even coordinates $u=(u_0,\dots,u_N)$.
Its full coefficients are $c_0=u_0$ and $c_{\pm i}=u_i/\sqrt2$; it has
boundary sum $S=\sum_jc_j=u_0+\sqrt2\sum_{i=1}^Nu_i$. For $n>N$, $Ac$ is
even, so its component along $e_n$ is
$R_n:=\ip{e_n}{Ac}=\sqrt2(Ac)_n$, and its odd components vanish. Hence
$\norm{\Pi_{\abs j>J}Ac}_{\ell^2(\Z)}=\bigl(\sum_{n>J}R_n^2\bigr)^{1/2}$.

\begin{lemma}[Paired identity]\label{lem:paired}
For every $n>N$,
\begin{equation}\label{eq:paired}
R_n=\frac{\sqrt2\,b_nS}n+\frac{2b_n}n\sum_{i=1}^N\frac{u_i\,i^2}{n^2-i^2}-2\sum_{i=1}^N\frac{i\,b_iu_i}{n^2-i^2}.
\end{equation}
\end{lemma}

\begin{proof}
By~\eqref{eq:displacement},
$(Ac)_n=b_n\sum_jc_j/(n-j)-\sum_jb_jc_j/(n-j)$. Pair $j=\pm i$:
$\frac1{n-i}+\frac1{n+i}=\frac{2n}{n^2-i^2}$; and, since $b_{-i}=-b_i$,
$\frac{b_i}{n-i}-\frac{b_i}{n+i}=\frac{2ib_i}{n^2-i^2}$. The term $j=0$
contributes $u_0/n$ to the first sum and nothing to the second, because
$b_0=0$. Finally write $\frac{2n}{n^2-i^2}=\frac2n+\frac{2i^2}{n(n^2-i^2)}$ and
collect $\frac{\sqrt2}{n}(u_0+\sqrt2\sum u_i)=\frac{\sqrt2S}n$.
\end{proof}

\begin{definition}[Signed and absolute moments]
For an order $R\ge0$, set
$m_r=\sum_{i=1}^Nu_i\,i^{2r}$ for $1\le r\le R$, and
$\nu_r=\sum_{i=1}^Nu_ib_i\,i^{2r+1}$ for $0\le r\le R-1$. Further set
$M=\sum_{i=1}^N\abs{u_i}\,i^{2R+2}$ and $V=\sum_{i=1}^N\abs{u_ib_i}\,i^{2R+1}$.
\end{definition}

\begin{lemma}[Finite geometric expansion]\label{lem:geometric}
For $n>N$,
\[
R_n=\frac{\sqrt2b_nS}n+2b_n\sum_{r=1}^R\frac{m_r}{n^{2r+1}}-2\sum_{r=0}^{R-1}\frac{\nu_r}{n^{2r+2}}+\mathcal E_n,
\]
\[
\mathcal E_n=\frac{2b_n}{n^{2R+1}}\sum_{i=1}^N\frac{u_i\,i^{2R+2}}{n^2-i^2}-\frac2{n^{2R}}\sum_{i=1}^N\frac{u_ib_i\,i^{2R+1}}{n^2-i^2}.
\]
If $J>N$ and $\abs{b_n}\le B_J$ for all $n>J$, then for all $n>J$
\[
\abs{\mathcal E_n}\le2\gamma_J\Bigl[\frac{B_JM}{n^{2R+3}}+\frac V{n^{2R+2}}\Bigr],\qquad\gamma_J=\frac1{1-(N/J)^2}.
\]
\end{lemma}

\begin{proof}
With $x=i^2/n^2<1$ one has $\frac x{1-x}=\sum_{r=1}^Rx^r+\frac{x^{R+1}}{1-x}$ and
$\frac1{1-x}=\sum_{r=0}^{R-1}x^r+\frac{x^R}{1-x}$. Hence
$\frac{i^2}{n^2-i^2}=\sum_{r=1}^R\frac{i^{2r}}{n^{2r}}+\frac{i^{2R+2}}{n^{2R}(n^2-i^2)}$
and
$\frac1{n^2-i^2}=\sum_{r=0}^{R-1}\frac{i^{2r}}{n^{2r+2}}+\frac{i^{2R}}{n^{2R}(n^2-i^2)}$.
Substitute these into~\eqref{eq:paired}. For the bound, note that $i\le N<J<n$
gives $\frac1{n^2-i^2}=\frac{n^{-2}}{1-(i/n)^2}\le\frac{\gamma_J}{n^2}$.
\end{proof}

\begin{lemma}[Integral test]\label{lem:integral}
For integers $p\ge1$ and $J\ge1$,
$H_p(J):=\bigl(\sum_{n>J}n^{-2p}\bigr)^{1/2}\le J^{1/2-p}/\sqrt{2p-1}$.
\end{lemma}

\begin{proof}
$t\mapsto t^{-2p}$ decreases, so $\sum_{n>J}n^{-2p}\le\int_J^\infty t^{-2p}dt$.
\end{proof}

\begin{theorem}[Infinite-index tail bound]\label{thm:tail}
Let $N\ge1$, $R\ge0$, $J>N$, and $\abs{b_n}\le B_J$ for $n>J$. For every
real even $u\in\R^{N+1}$,
\begin{multline}\label{eq:tail}
\Bigl(\sum_{n>J}R_n^2\Bigr)^{1/2}\le\sqrt2B_J\abs SH_1+2B_J\sum_{r=1}^R\abs{m_r}H_{2r+1}\\
+2\sum_{r=0}^{R-1}\abs{\nu_r}H_{2r+2}+2\gamma_J\bigl[B_JMH_{2R+3}+VH_{2R+2}\bigr],
\end{multline}
with $H_p=H_p(J)$.
\end{theorem}

\begin{proof}
Apply Minkowski's inequality in $\ell^2(\{n>J\})$ to the decomposition of
Lemma~\ref{lem:geometric}. Bound each sequence pointwise by a constant times
$n^{-p}$; then Lemma~\ref{lem:integral} gives its $\ell^2$ norm.
\end{proof}

Theorem~\ref{thm:tail} is a statement about every omitted index at once.
It holds for any finite rank, any cutoff $J>N$ and any order $R$. Neither
the sum over $n$ nor the remainder is truncated. The only inputs it needs
are the scalars $S$, $m_r$, $\nu_r$, $M$, $V$ and $B_J$.

\begin{remark}[Why the signs matter]\label{rem:uncoupled}
Bounding $\abs{A_{nj}}\le2B/\abs{n-j}$ entrywise and summing gives the
\emph{uncoupled} estimate
$\bigl(\sum_{n>J}R_n^2\bigr)^{1/2}\le2\sqrt2B\bigl(\abs{u_0}+\sqrt2\sum\abs{u_i}\bigr)/\sqrt{J-N}$.
For the rank-eight ground on the cell this is about $2.906$ when $J=64$.
Theorem~\ref{thm:tail} gives $9.49\times10^{-12}$
(Theorem~\ref{thm:main-tail}), roughly $3.06\times10^{11}$ times smaller. The
improvement comes from three facts: the ground state has a tiny boundary
sum $S$, low signed moments $m_1$ and $\nu_0$ that are small, and the
moment expansion preserves the sign cancellation across the retained
coordinates.
\end{remark}

%======================================================================
\section{High-mode coercivity}\label{sec:coercive}
%======================================================================

\begin{lemma}[Pole form]\label{lem:pole-form}
For $f\in\mathcal H_L$, let $\tilde f(x)=f(x+L/2)$ on $[-L/2,L/2]$. Then
$\ip{c}{Pc}=2\abs{\ip{\cosh\frac\cdot2}{\tilde f}}^2-2\abs{\ip{\sinh\frac\cdot2}{\tilde f}}^2$.
Consequently $P\ge0$ on even vectors, and $P\ge-(2\sinh\frac L2-L)$ on
all vectors.
\end{lemma}

\begin{proof}
By~\eqref{eq:W02}, in additive coordinates,
$W_{0,2}(f^{*}*f)=\iint\overline{f(s)}f(t)\,2\cosh\frac{s-t}2\,ds\,dt$. This
kernel is translation invariant, so we may centre. Then
$2\cosh\frac{s-t}2=2\cosh\frac s2\cosh\frac t2-2\sinh\frac s2\sinh\frac t2$.
The two functions are orthogonal on $[-L/2,L/2]$. For even $\tilde f$ the
sinh term vanishes. In general the negative eigenvalue is
$-2\norm{\sinh\frac\cdot2}^2=-(2\sinh\frac L2-L)$.
\end{proof}

\begin{lemma}[Prime form]\label{lem:prime-form}
$\abs{\ip c{Sc}}\le2\mathcal P\norm c^2$ on the cell.
\end{lemma}

\begin{proof}
$\ip c{Sc}=\sum_{p^a<C}\frac{\log p}{p^{a/2}}\bigl(F(a\log p)+F(-a\log p)\bigr)$
with $F(y)=\int\overline{f(x-y)}f(x)\,dx$, and Cauchy--Schwarz gives
$\abs{F(y)}\le\norm f^2$.
\end{proof}

\begin{lemma}[Archimedean diagonal]\label{lem:arch-diag}
For $n\ne0$ and $C\in[17,19]$,
\[
-H_{nn}\ge\delta(\abs n):=\log\frac{\abs n}{\Lmax}-\frac{\Lmax}{\pi\abs n}-\frac1{4\abs n}-\frac{\Lmax(1+E)}{2\pi^2n^2}.
\]
\end{lemma}

\begin{proof}
Use Proposition~\ref{prop:arch} with $y=\pi\abs n/L$. By
Lemma~\ref{lem:digamma}(b), $-\log\pi+\Rea\psi(z_n)\ge\log\frac{\abs n}L-\frac L{\pi\abs n}$.
By Lemma~\ref{lem:digamma}(c),
$\frac{\abs{\psi'(z_n)}}{2L}\le\frac L{2\pi^2n^2}+\frac1{4\abs n}$. The
exponential series is at most
$\frac2L\sum_ke^{-a_kL}w_n^{-2}\le\frac{LE}{2\pi^2n^2}$, because
$\abs{a^2-w^2}\le a^2+w^2$. Each resulting term is monotone in
$L\in[\Lmin,\Lmax]$ in the direction used.
\end{proof}

\begin{theorem}[High-mode coercivity]\label{thm:coercive}
For $n_0\ge1$ put
\[
\beta_{\mathrm{even}}(n_0)=\delta(n_0)-2\pi B_\R-2\mathcal P,\qquad
\beta_{\mathrm{full}}(n_0)=\beta_{\mathrm{even}}(n_0)-\bigl(2\sinh\tfrac\Lmax2-\Lmax\bigr).
\]
For every $C\in[17,19]$ and every finitely supported $x$ with $x_j=0$ for
$\abs j<n_0$, one has $\ip x{A(C)x}\ge\beta_{\mathrm{full}}(n_0)\norm x^2$.
If moreover $x$ is even, then $\ip x{A(C)x}\ge\beta_{\mathrm{even}}(n_0)\norm x^2$.
Both lower bounds increase to $+\infty$ as $n_0$ grows, and
\[
\beta_{\mathrm{even}}(18{,}612{,}929)>1.000000028,\qquad
\beta_{\mathrm{full}}(60{,}879{,}750)>1.000000002.
\]
\end{theorem}

\begin{proof}
Decompose $\ip xAx=\ip xPx-\ip x{H_{\mathrm{diag}}x}-\ip x{H_{\mathrm{off}}x}-\ip xSx$.
Then apply Lemma~\ref{lem:pole-form}, Lemma~\ref{lem:arch-diag} (with $\delta$
increasing in $\abs n$), Corollary~\ref{cor:offdiag} (compressions keep the
bound) and Lemma~\ref{lem:prime-form}. The numerical values come from
outward interval evaluation (Table~\ref{tab:constants}). The certificate
uses 1024-bit Arb; the independent check uses mpmath intervals. Both report
$\beta_{\mathrm{even}}(18{,}612{,}928)<1$ and $\beta_{\mathrm{full}}(60{,}879{,}749)<1$.
Thus the stated starts are the least integers for which \emph{these
particular} bounds exceed one; they are not claimed sharp for $A$ itself.
\end{proof}

Since the span of the $U_n$ is a form core for the closed semilocal
form~\cite[Prop.~3.4]{CCM2025}, the same lower bounds hold for the closure
of the far-block form. We use only finite compressions below.

%======================================================================
\section{The finite-rank certificate}\label{sec:certificate}
%======================================================================

Theorem~\ref{thm:tail} reduces the infinite tail to finitely many real
numbers attached to $u(C)$. This section proves that the enclosures of
those numbers are valid uniformly in $C\in[17,19]$. All arithmetic is
midpoint--radius ball arithmetic with outward rounding (Arb via
python-flint 0.8.0, 1024 bits~\cite{Johansson2017}). The
\emph{inclusion property} is the only fact about the computer we use: each
computed ball contains the exact value of the corresponding real
expression for every choice of inputs from the input balls.

\subsection{Analytic Taylor enclosure in the support parameter}
Replace $\Rea\psi(z_n)$ by $\frac12[\psi(\frac14-\frac{i\pi n}L)+\psi(\frac14+\frac{i\pi n}L)]$,
and $\Ima\psi$ by the corresponding odd combination divided by $2i$. Then
every entry of Section~\ref{sec:closed} is holomorphic in $L$ on $\Rea L>0$,
for a fixed active prime set. Let $c_\ast=\frac12(\Lmin+\Lmax)$,
$\delta=\frac12(\Lmax-\Lmin)<0.055613$ and $R_\ast=\frac14$.

\begin{lemma}[Cauchy remainder]\label{lem:cauchy}
Let $W$ be holomorphic on $\abs{L-c_\ast}\le R_\ast$ with $\abs W\le M_W$
there, and let $W_j$ be its Taylor coefficients at $c_\ast$. If
$\abs h\le\delta<R_\ast$ and $q=\delta/R_\ast$, then
$\abs{W(c_\ast+h)-\sum_{j\le d}W_jh^j}\le M_Wq^{d+1}/(1-q)$.
\end{lemma}

\begin{proof}
Cauchy's estimate gives $\abs{W_j}\le M_W/R_\ast^j$; sum the geometric tail.
\end{proof}

The certificate~\cite[artifact \texttt{weil\_support\_taylor\_17\_19\_8\_degree64}]{Artifacts}
encloses every coefficient of the $9\times9$ even block to degree $d=64$.
It truncates the exponential series after 80 terms, with an omitted tail
bounded uniformly over the complex disk and over \emph{all} index pairs. It
bounds $M_W$ on the circle by covering it with rectangles. The resulting
operator-norm error of the polynomial matrix over the real cell is
below $4.870\times10^{-41}$.

\subsection{Validated eigenfamily}
The cell is covered by 32 adjacent closed subcells. On each subcell the
polynomial matrix is translated to the local centre. A candidate eigenbasis
$U(L)$ and candidate eigenvalues $\Lambda=\diag(\ell_i(L))$ are polynomials
with exact dyadic coefficients, written in a fixed orthogonal centre basis
$Q$. No convergence of the recurrence that proposed them is assumed.

\begin{lemma}[Uniform isolation]\label{lem:isolation}
Let $B=Q^TAQ$ be real symmetric. Let $E=BU-U\Lambda$, and suppose
$\norm{U-I}_\infty\le\beta<1$. Put $\eta=\norm E_\infty/(1-\beta)$. If the
intervals $[\ell_i-\eta,\ell_i+\eta]$ are pairwise disjoint, then each
contains exactly one eigenvalue of $A$, in increasing order.
\end{lemma}

\begin{proof}
By the Neumann series, $\norm{U^{-1}}_\infty\le1/(1-\beta)$, and
$U^{-1}BU=\Lambda+U^{-1}E$. Gershgorin's theorem with row sums places the
spectrum in the union of the disks of radius $\eta$ about the $\ell_i$. By
continuity of eigenvalues along $t\mapsto\Lambda+tU^{-1}E$, a union of $k$
disjoint disks contains exactly $k$ eigenvalues. The eigenvalues are real
because $A$ is symmetric.
\end{proof}

\begin{lemma}[Eigenvector distance]\label{lem:distance}
Let $A$ be real symmetric with simple lowest eigenvalue $\lambda_0$ and unit
eigenvector $v$. Let $\hat u$ be a unit vector with
$\norm{A\hat u-\ell\hat u}\le\rho$, and suppose every other eigenvalue
satisfies $\lambda-\ell\ge D>0$. If $v$ is oriented so that
$\ip v{\hat u}\ge0$, then $\norm{\hat u-v}\le\sqrt2\rho/D$.
\end{lemma}

\begin{proof}
Write $\hat u=\alpha v+w$ with $w\perp v$ and $\alpha\ge0$. Projecting
$(A-\ell)\hat u$ onto $v^\perp$ gives $(A-\ell)w$, so $D\norm w\le\rho$.
Then $\norm{\hat u-v}^2=2(1-\alpha)\le2(1-\alpha^2)=2\norm w^2$.
\end{proof}

All 32 subcells pass Lemma~\ref{lem:isolation}, with
$\norm{U-I}_\infty<0.01633$ and disk radii below $4.456\times10^{-40}$. An
independent odd-block certificate with the same subdivision places every
odd eigenvalue above the even ground. Hence $u(C)$ is the simple ground of
the full $17\times17$ truncation; the tail theorem below does not need this.
Uniformly on the cell, the lowest eigenvalue lies in
$[2.259\times10^{-26},\,1.380\times10^{-24}]$, and the full ground gap
exceeds $4.043\times10^{-23}$. Lemma~\ref{lem:distance} gives
\begin{equation}\label{eq:dist}
\norm{\hat u(C)-u(C)}\le d<8.687\times10^{-21},
\end{equation}
where $\hat u=p/\norm p$ is the normalized candidate. The candidate's
centre-basis anchor coordinate is identically one, so
$\norm{p}\in[1,\bar\nu]$ for a certified $\bar\nu$. The unit ground origin
satisfies $u_0>0.5753$, which fixes the orientation consistently across
subcells.

\subsection{Enclosures of the moments}

\begin{lemma}[Linear functionals of the exact ground]\label{lem:functional}
Let $\omega\in\R^{N+1}$. For $C$ in a subcell,
\[
\ip\omega{u(C)}\in\frac{\mathrm{range}\,\ip\omega{p}}{[1,\bar\nu]}+\bigl[-d\norm\omega,\,d\norm\omega\bigr].
\]
If the weights depend on $C$ through $b_i(L)=\mathrm{poly}_i(L)+\eta_i$ with
$\abs{\eta_i}\le\epsilon_i$, then
\[
\sum_i\omega_ib_iu_i\in\frac{\mathrm{range}\sum_i\omega_i\,\mathrm{poly}_i\,p_i}{[1,\bar\nu]}+\Bigl[\pm\Bigl(d\,\bigl\|(\omega_i\sup\abs{b_i})_i\bigr\|+\sum_i\abs{\omega_i}\epsilon_i\sup\abs{p_i}\Bigr)\Bigr].
\]
\end{lemma}

\begin{proof}
Write $u=\hat u+(u-\hat u)$, use~\eqref{eq:dist} with Cauchy--Schwarz, and
use $\norm p\ge1$ for the polynomial-error term.
\end{proof}

The key design choice is that each moment is formed as \emph{one}
polynomial in the support parameter, namely $\sum_i\omega_ip_i(L)$ or
$\sum_i\omega_ib_i(L)p_i(L)$, \emph{before} its range is enclosed. Neighbouring
Fourier coordinates nearly cancel, and a common parameter is shared, so
enclosing each coordinate separately would destroy the cancellation that
Remark~\ref{rem:uncoupled} quantifies. We apply the lemma with:
\begin{itemize}
\item $\omega=(1,\sqrt2,\dots,\sqrt2)$ for $S$. The spatial certificate
evaluates $S$ as the endpoint Bernstein coefficient of the spatial ground
kernel, which equals $S$ exactly;
\item $\omega_i=i^{2r}$ ($i\ge1$) for $m_r$, $1\le r\le8$;
\item $\omega_i=i^{2r+1}$ with $b_i(L)=iA_{i0}(L)$ taken from the same
Taylor certificate, for $\nu_r$, $0\le r\le7$;
\item absolute weights after the signed expansion, for $M$ and $V$ with $R=8$.
\end{itemize}

\subsection{The main computer-assisted theorem}

\begin{theorem}[Coupled infinite Fourier tail]\label{thm:main-tail}
For every $C\in[17,19]$, the unit even rank-eight ground vector $u(C)$
satisfies
\[
\Bigl(\sum_{n>64}R_n(C)^2\Bigr)^{1/2}\le 9.4900466\times10^{-12}<9.491\times10^{-12}.
\]
\end{theorem}

\begin{proof}
On each of the 32 subcells, apply Theorem~\ref{thm:tail} with $N=8$, $J=64$,
$R=8$ and $B_{64}=B_0+B_1/64$; this is valid because $n>64$ implies
$\abs{b_n}\le B_0+B_1/n<B_{64}$ by Proposition~\ref{prop:b-bounds}. The scalar
inputs are enclosed uniformly over the subcell by Lemma~\ref{lem:functional}.
The right side of~\eqref{eq:tail} is evaluated in outward ball arithmetic,
and the bound is the maximum over the subcells. Table~\ref{tab:worst-cell}
lists the contributions for the worst subcell, which is adjacent to $C=17$.
\end{proof}

\begin{table}[ht]
\centering
\caption{Worst subcell (index 0, adjacent to $C=17$) of
Theorem~\ref{thm:main-tail}. Each entry is a certified outward upper bound.
The moment intervals are certified enclosures over the whole subcell.}
\label{tab:worst-cell}
\begin{tabular}{@{}lll@{}}
\toprule
Term of~\eqref{eq:tail} & Certified input & Contribution\\
\midrule
$\sqrt2B_{64}\abs SH_1$ & $S\in[6.3167,\,6.5920]\times10^{-12}$ & $2.465\times10^{-12}$\\
$2B_{64}\abs{m_1}H_3$ & $m_1\in[-6.317,-5.976]\times10^{-8}$ & $3.648\times10^{-12}$\\
$2B_{64}\abs{m_2}H_5$ & $m_2\in[1.1354,\,1.1920]\times10^{-4}$ & $1.252\times10^{-12}$\\
$2B_{64}\abs{m_3}H_7$ & $m_3\in[-7.691,-7.419]\times10^{-2}$ & $1.642\times10^{-13}$\\
$2B_{64}\sum_{r=4}^8\abs{m_r}H_{2r+1}$ & & $<9.5\times10^{-15}$\\
$2\abs{\nu_0}H_2$ & $\nu_0\in[-5.406,-4.750]\times10^{-10}$ & $1.219\times10^{-12}$\\
$2\abs{\nu_1}H_4$ & $\nu_1\in[1.5495,\,1.7150]\times10^{-6}$ & $6.182\times10^{-13}$\\
$2\abs{\nu_2}H_6$ & $\nu_2\in[-1.5124,-1.3961]\times10^{-3}$ & $1.062\times10^{-13}$\\
$2\sum_{r=3}^7\abs{\nu_r}H_{2r+2}$ & & $<7.9\times10^{-15}$\\
remainder $2\gamma_{64}[\cdots]$ & $M,V$ & $1.001\times10^{-22}$\\
\midrule
total & & $9.4900\times10^{-12}$\\
\bottomrule
\end{tabular}
\end{table}

Over all 32 subcells the bound ranges from $3.03\times10^{-12}$ to
$9.490\times10^{-12}$, and the expansion remainder never exceeds
$1.219\times10^{-22}$. The uncoupled estimate of
Remark~\ref{rem:uncoupled} ranges from $2.889$ to $2.906$.

%======================================================================
\section{A consequence, and what remains open}\label{sec:limits}
%======================================================================

\begin{corollary}[Far-block inverse energy]\label{cor:energy}
Let $n_0=18{,}612{,}929$. Let $Q_M$ be the orthogonal projection onto
$\mathrm{span}\{e_j:n_0\le j\le M\}$ for any $M\ge n_0$, let $\lambda(C)$ be
the lowest even eigenvalue at rank eight, and let
$T_M=Q_MAQ_M-\lambda I$ on the range of $Q_M$. Then, uniformly in $M$ and
$C\in[17,19]$,
\[
\ip{Q_MAc}{T_M^{-1}Q_MAc}\le\frac{\norm{Q_MAc}^2}{1-\lambda}<9.007\times10^{-23}.
\]
\end{corollary}

\begin{proof}
By Theorem~\ref{thm:coercive}, $Q_MAQ_M\ge\beta_{\mathrm{even}}(n_0)>1$, so
$T_M\ge1-\lambda$, with $0<\lambda<1.380\times10^{-24}$. Since $n_0>64$,
Theorem~\ref{thm:main-tail} gives $\norm{Q_MAc}^2<(9.49005\times10^{-12})^2<9.0061\times10^{-23}$.
\end{proof}

This is the energy of the \emph{direct} far-block response. It is not the
change of the ground state when the rank is increased. The following
problems remain open.

\begin{problem}[Intermediate block]\label{prob:intermediate}
Control the Schur complement obtained by eliminating the modes
$9\le j<18{,}612{,}929$, together with the profile response it induces. The
forcing into those modes is not bounded here; for $9\le j\le64$ it is a
finite computation, but for $j>64$ only its $\ell^2$ norm is controlled.
More importantly, the eliminated block contains spectrum close to the
ground eigenvalue. Tiny energy does not imply a tiny profile increment;
the response can be amplified by a nearly singular inverse. A
Davis--Kahan bound~\cite{DavisKahan1970} that uses the full residual would
need a spectral gap of the infinite operator larger than the forcing,
whereas the certified finite gap is about $4\times10^{-23}$.
\end{problem}

\begin{problem}[All supports and ranks]
Find estimates of Theorem~\ref{thm:main-tail} type that are uniform as
$C\to\infty$ and $N\to\infty$ jointly, with a summable combined
support/rank path. The present constants $\mathcal P$, $D_{+}$ and the
coercivity start $n_0$ all grow with $C$, so the argument does not
extend without new ideas.
\end{problem}

\begin{problem}[Identification with $\Xi$]\label{prob:xi}
Show that suitably normalized ground profiles converge to Riemann's
$\Xi$ function, as proposed in~\cite[\S7]{CCM2025}. Nothing in this paper
addresses that identification.
\end{problem}

Consequently, the results here are finite-cell, rank-independent
estimates. They are not, and do not imply, a proof of the Riemann
Hypothesis.

\subsection*{Theorem-status ledger}
\begin{center}
\begin{tabular}{@{}p{0.47\textwidth}p{0.45\textwidth}@{}}
\toprule
Statement & Status\\
\midrule
Lemmas~\ref{lem:halfline}--\ref{lem:corr}, Props.~\ref{prop:displacement}--\ref{prop:primes}, Cor.~\ref{cor:b} & Proved (from the explicit-formula distributions)\\
Lemma~\ref{lem:digamma}, Prop.~\ref{prop:b-bounds}, Thm.~\ref{thm:hilbert} & Proved; constants by interval evaluation\\
Lemmas~\ref{lem:paired}--\ref{lem:integral}, Thm.~\ref{thm:tail} & Proved, all ranks, all indices\\
Thm.~\ref{thm:coercive} & Proved; thresholds by interval evaluation\\
Lemmas~\ref{lem:cauchy}--\ref{lem:functional} & Proved\\
Thm.~\ref{thm:main-tail}, Cor.~\ref{cor:energy} & Proved, computer-assisted (validated ball arithmetic)\\
Appendix~\ref{app:crosscheck} & Independent floating-point evidence, not proof\\
Problems~\ref{prob:intermediate}--\ref{prob:xi}, RH & Open\\
\bottomrule
\end{tabular}
\end{center}

\subsection*{Software and use of AI tools}
Ball arithmetic uses Arb through python-flint~0.8.0~\cite{Johansson2017};
the independent cross-check uses mpmath. The author used AI tools to write
and execute the code for the certificate generators, the regression tests
and the independent cross-check. Every computational claim in this paper
is reproducible from the replay commands in Appendix~\ref{app:repro}.

\appendix

%======================================================================
\section{Independent cross-check}\label{app:crosscheck}
%======================================================================

The script \texttt{paper/weil\_tail\_crosscheck.py} shares no code with the
certificate generator. It implements the formulas of
Sections~\ref{sec:closed}--\ref{sec:tail} directly in mpmath at 80 digits.
Its results, which are floating-point evidence and not part of any proof:
\begin{enumerate}
\item \emph{Closed entries versus the defining distributions.} At $C=17$
the closed entries $(0,0)$, $(3,3)$, $(8,8)$, $(1,0)$, $(5,-2)$, $(8,-8)$ and
$(7,6)$ agree with adaptive quadrature of~\eqref{eq:Dist}--\eqref{eq:Arch} to
within $1.6\times10^{-30}$. This checks the $\log\pi$ regularization of
Proposition~\ref{prop:arch} independently.
\item \emph{Paired identity.} At $C=17,18,19$ and $n=9,65,1000$,
\eqref{eq:paired} agrees with a direct product $\sqrt2\sum_jA_{nj}c_j$ to
within $10^{-60}$.
\item \emph{Ground moments.} The independently computed unit ground at
$C=17$, $18$ and $19$ gives
{\small\[
\setlength{\arraycolsep}{4pt}
\begin{array}{@{}lccccc@{}}
C & S & m_1 & \nu_0 & m_2 & \nu_1\\
17 & 6.3425\times10^{-12} & -6.2547\times10^{-8} & -5.3873\times10^{-10} & 1.1860\times10^{-4} & 1.7116\times10^{-6}\\
18 & 1.3773\times10^{-12} & -1.9552\times10^{-8} & -1.2350\times10^{-10} & 5.2081\times10^{-5} & 5.3834\times10^{-7}\\
19 & 4.0274\times10^{-12} & -4.1936\times10^{-8} & -3.4992\times10^{-10} & 8.3493\times10^{-5} & 1.1650\times10^{-6}
\end{array}
\]}
In each case $S$, $m_1$ and $\nu_0$ lie inside the certified intervals of the
containing subcell (indices 0, 16 and 31).
\item \emph{Direct omitted-mode sum.} Summing $R_n^2$ directly for
$64<n\le20{,}000$ and bounding $n>20{,}000$ by Theorem~\ref{thm:tail} with
$R=4$ gives totals $\le1.02\times10^{-12}$, $2.02\times10^{-13}$ and
$6.45\times10^{-13}$ at $C=17,18,19$. These are consistent with the
certified uniform bound, and the direct sum at $C=17$
($8.21\times10^{-13}$) shows that the bound is conservative by a factor of
about $11$.
\item \emph{Constants.} The mpmath interval evaluation reproduces every entry
of Table~\ref{tab:constants}. It also reproduces both coercivity thresholds,
including the failure at $n_0-1$.
\end{enumerate}

%======================================================================
\section{Replay and data}\label{app:repro}
%======================================================================

The certificate chain (matrix Taylor, even eigenfamily, odd block, spatial
boundary, rank tail) is stored in \texttt{research/prime\_spectral/}. Each
artifact records the SHA-256 hashes of its source files and of its input
artifacts, and the generators refuse stale inputs. The tail certificate
\texttt{weil\_rank\_tail\_cell\_17\_19\_8\_degree64.json} records the input
hashes
\begin{center}\small
\begin{tabular}{@{}ll@{}}
Taylor & \texttt{5f745aa6175bc5b06477d90e587f2435b74feae4a1012af70710bacaa573e0db}\\
even & \texttt{8a977827a832ccaca8c706eab1ad334e4636d2508518b21f3e708a6a39a1083b}\\
\end{tabular}
\end{center}
Replay the tail certificate (7 tests, including a full replay of the 32-cell
cover) and the independent check with
\begin{verbatim}
uv run --with python-flint==0.8.0 --python 3.12 python -m unittest \
  discover -s experiments/prime_spectral -p test_weil_rank_tail_cell.py -v
uv run --with mpmath python paper/weil_tail_crosscheck.py 20000
\end{verbatim}
The full prime-spectral regression suite (159 tests) is run with the same
command and the pattern \texttt{test\_*.py}. The analytic arguments of this
paper are not formalized in Lean.

\begin{thebibliography}{99}
\bibitem{Artifacts}
B.~Mondal, \emph{Reinmann: certificates and replay scripts for the semilocal Weil form}, research repository, directories \texttt{experiments/prime\_spectral} (code and certificates) and \texttt{research/prime\_spectral} (reports), 2026. Archived with this preprint, doi:10.5281/zenodo.23187645.

\bibitem{Bombieri2000}
E.~Bombieri, Remarks on Weil's quadratic functional in the theory of prime numbers.~I, \emph{Rend. Mat. Acc. Lincei} (9) \textbf{11} (2000), 183--233.

\bibitem{CCM2025}
A.~Connes, C.~Consani and H.~Moscovici, Zeta spectral triples, \emph{arXiv:2511.22755} (2025).

\bibitem{DavisKahan1970}
C.~Davis and W.~M.~Kahan, The rotation of eigenvectors by a perturbation.~III, \emph{SIAM J. Numer. Anal.} \textbf{7} (1970), 1--46.

\bibitem{DLMF}
\emph{NIST Digital Library of Mathematical Functions}, \url{https://dlmf.nist.gov/}, \S\S5.7, 5.9, 5.15, 6.6.

\bibitem{HLP}
G.~H.~Hardy, J.~E.~Littlewood and G.~P\'olya, \emph{Inequalities}, 2nd ed., Cambridge University Press, 1952.

\bibitem{Johansson2017}
F.~Johansson, Arb: efficient arbitrary-precision midpoint-radius interval arithmetic, \emph{IEEE Trans. Comput.} \textbf{66} (2017), 1281--1292.

\bibitem{Rump2010}
S.~M.~Rump, Verification methods: rigorous results using floating-point arithmetic, \emph{Acta Numer.} \textbf{19} (2010), 287--449.

\bibitem{Weil1952}
A.~Weil, Sur les ``formules explicites'' de la th\'eorie des nombres premiers, \emph{Comm. S\'em. Math. Univ. Lund} (1952), tome suppl\'ementaire, 252--265.
\end{thebibliography}

\end{document}
